% Options for packages loaded elsewhere
\PassOptionsToPackage{unicode}{hyperref}
\PassOptionsToPackage{hyphens}{url}
\PassOptionsToPackage{table}{xcolor}
%
\documentclass[
  11pt,
]{report}
\usepackage{amsmath,amssymb}
\usepackage{iftex}
\ifPDFTeX
  \usepackage[T1]{fontenc}
  \usepackage[utf8]{inputenc}
  \usepackage{textcomp} % provide euro and other symbols
\else % if luatex or xetex
  % pdfLaTeX-only repository policy: the original unicode-math/fontspec
  % branch is unreachable under pdfLaTeX and is replaced with a guard.
  \PackageError{cpython-suite}{Compile with pdfLaTeX only}{}
\fi
\usepackage{lmodern}
\ifPDFTeX\else
  % xetex/luatex font selection
\fi
% Use upquote if available, for straight quotes in verbatim environments
\IfFileExists{upquote.sty}{\usepackage{upquote}}{}
\IfFileExists{microtype.sty}{% use microtype if available
  \usepackage[]{microtype}
  \UseMicrotypeSet[protrusion]{basicmath} % disable protrusion for tt fonts
}{}
\makeatletter
\@ifundefined{KOMAClassName}{% if non-KOMA class
  \IfFileExists{parskip.sty}{%
    \usepackage{parskip}
  }{% else
    \setlength{\parindent}{0pt}
    \setlength{\parskip}{6pt plus 2pt minus 1pt}}
}{% if KOMA class
  \KOMAoptions{parskip=half}}
\makeatother
\usepackage{xcolor}
\usepackage{listings}
\newcommand{\passthrough}[1]{#1}
\lstset{defaultdialect=[5.3]Lua}
\lstset{defaultdialect=[x86masm]Assembler}
\usepackage{longtable,booktabs,array}
\usepackage{calc} % for calculating minipage widths
% Correct order of tables after \paragraph or \subparagraph
\usepackage{etoolbox}
\makeatletter
\patchcmd\longtable{\par}{\if@noskipsec\mbox{}\fi\par}{}{}
\makeatother
% Allow footnotes in longtable head/foot
\IfFileExists{footnotehyper.sty}{\usepackage{footnotehyper}}{\usepackage{footnote}}
\makesavenoteenv{longtable}
\setlength{\emergencystretch}{3em} % prevent overfull lines
\providecommand{\tightlist}{%
  \setlength{\itemsep}{0pt}\setlength{\parskip}{0pt}}
\setcounter{secnumdepth}{5}
% CPython Engineering Suite shared presentation layer -- pdfLaTeX compatible
\usepackage[T1]{fontenc}
\usepackage{lmodern}
\usepackage{microtype}
\usepackage{geometry}
\geometry{letterpaper,margin=0.91in,headheight=16pt,headsep=19pt}
\usepackage{xcolor} % [table] passed at top of file (option-clash fix)
\definecolor{CPNavy}{HTML}{13283B}
\definecolor{CPBlue}{HTML}{176BA0}
\definecolor{CPGold}{HTML}{D8A73C}
\definecolor{CPMuted}{HTML}{566979}
\definecolor{CPBg}{HTML}{EFF5F8}
\usepackage{titlesec}
\titleformat{\chapter}[display]{\normalfont\sffamily\bfseries\color{CPNavy}}{\large\color{CPBlue}CHAPTER \thechapter}{9pt}{\Huge}[\vspace{6pt}\color{CPGold}\titlerule]
\titleformat{\section}{\Large\bfseries\sffamily\color{CPBlue}}{\thesection}{0.7em}{}
\titleformat{\subsection}{\large\bfseries\sffamily\color{CPNavy}}{\thesubsection}{0.7em}{}
\usepackage{fancyhdr}
\pagestyle{fancy}
\fancyhf{}
\fancyhead[L]{\footnotesize\color{CPMuted}\sffamily CPYTHON ENGINEERING SUITE}
\fancyhead[R]{\footnotesize\color{CPMuted}\sffamily 3.16.0a1+dev SNAPSHOT}
\fancyfoot[C]{\footnotesize\color{CPMuted}\thepage}
\renewcommand{\headrulewidth}{0.3pt}
\AtBeginDocument{\urlstyle{tt}}
\lstset{basicstyle=\ttfamily\small,columns=fullflexible,keepspaces=true,breaklines=true,breakatwhitespace=false,showstringspaces=false,frame=single,framesep=5pt,rulecolor=\color{CPBg},backgroundcolor=\color{CPBg},xleftmargin=3pt,xrightmargin=3pt}
\usepackage{tcolorbox}
\tcbuselibrary{breakable,skins}
\newtcolorbox{sourcebox}[1][]{enhanced,breakable,sharp corners,boxrule=0.4pt,colback=CPBg,colframe=CPBlue,borderline west={3pt}{0pt}{CPBlue},left=10pt,right=8pt,top=6pt,bottom=6pt,#1}
\newtcolorbox{cautionbox}[1][]{enhanced,breakable,sharp corners,boxrule=0.4pt,colback=CPBg,colframe=CPGold,borderline west={3pt}{0pt}{CPGold},left=10pt,right=8pt,top=6pt,bottom=6pt,#1}
\usepackage{tikz}
\usetikzlibrary{arrows.meta,positioning,shapes.geometric,calc}
\usepackage{enumitem}
\setlist{itemsep=2pt,topsep=4pt}
\usepackage{booktabs}
\usepackage{longtable}
\AtBeginEnvironment{longtable}{\small}
\usepackage{array}
\setlength{\parskip}{5pt}
\setlength{\parindent}{0pt}
\setcounter{tocdepth}{2}
\usepackage{hyperref}
\hypersetup{colorlinks=true,linkcolor=CPBlue,urlcolor=CPBlue,citecolor=CPBlue,pdfauthor={CPython Engineering Documentation Suite}}
\makeatletter
\renewcommand{\maketitle}{%
  \begin{titlepage}\pagecolor{CPNavy}\color{white}\sffamily
  \vspace*{0.8in}
  {\color{CPGold}\rule{2.1in}{4pt}}\\[16pt]
  {\Large\bfseries CPYTHON ENGINEERING DOCUMENTATION SUITE}\\[24pt]
  {\Huge\bfseries \@title\par}\vspace{19pt}
  \vfill
  {\normalsize Source baseline: CPython 3.16.0 alpha 1+dev repository snapshot}\par
  {\normalsize Prepared 16 October 2026}\par
  {\small\color{white!70} Source-traced engineering commentary; not an official Python specification.}\vspace{0.65in}
  \end{titlepage}\nopagecolor\color{black}
  }
\makeatother
\ifLuaTeX
  \usepackage{selnolig}  % disable illegal ligatures
\fi
\usepackage{bookmark}
\IfFileExists{xurl.sty}{\usepackage{xurl}}{} % add URL line breaks if available
\urlstyle{same}
\hypersetup{
  pdftitle={CPython Memory Management and Concurrency},
  hidelinks,
  pdfcreator={LaTeX via pandoc}}

\title{CPython Memory Management and Concurrency}
\usepackage{etoolbox}
\makeatletter
\providecommand{\subtitle}[1]{% add subtitle to \maketitle
  \apptocmd{\@title}{\par {\large #1 \par}}{}{}
}
\makeatother
\subtitle{Volume 05 \textbar{} Reference ownership, cyclic GC, GIL, free
threading and QSBR}
\author{}
\date{}

\begin{document}
\maketitle

{
\setcounter{tocdepth}{1}
\tableofcontents
}
\chapter{Behavior guarantees versus runtime
mechanics}\label{behavior-guarantees-versus-runtime-mechanics}

\begin{sourcebox}
\textbf{Evidence classification.} \textbf{SOURCE} = statements supported by material present in the uploaded snapshot; \textbf{IMPLEMENTATION} = CPython-specific internals, not guaranteed Python semantics; \textbf{ENGINEERING} = original recommendations, examples, checks, and exercises. The source baseline is \textbf{3.16.0 alpha 1+dev}; missing repository file bodies are noted as evidence limitations.
\end{sourcebox}

CPython's memory management combines reference counting, cyclic garbage
collection, allocators and, in free-threaded configurations, additional
concurrency-safe lifetime-management mechanisms. The details in
\path{InternalDocs/garbage_collector.md} and
\path{InternalDocs/qsbr.md} are
\textbf{implementation-specific}. Python-language programs should use
explicit resource-management protocols rather than assuming the timing
of object finalization.

\textbf{Snapshot:} CPython 3.16.0 alpha 1+dev. Compare the documented
build configuration and extension compatibility of the \emph{actual
deployed interpreter} before generalizing measurements or thread-safety
behavior.

\chapter{Reference counting and object
lifetime}\label{reference-counting-and-object-lifetime}

\path{InternalDocs/garbage_collector.md} describes
reference counting as the main collection mechanism: when the last
relevant ownership relationship disappears, deallocation can occur;
objects that themselves own references can cascade decrements. The C API
exposes reference-management operations such as
\passthrough{\lstinline!Py\_INCREF!},
\passthrough{\lstinline!Py\_DECREF!},
\passthrough{\lstinline!Py\_NewRef!},
\passthrough{\lstinline!Py\_XDECREF!} and ownership-transfer
conventions. The public docs explain \textbf{strong (owned)} versus
\textbf{borrowed} references; these determine whether the caller must
release a reference.

\begin{lstlisting}[language=C]
/* Illustrative ownership pattern; use in an extension with Python.h. */
static PyObject *identity_newref(PyObject *object)
{
    return Py_NewRef(object);  /* caller owns returned reference */
}
\end{lstlisting}

Do not use raw \passthrough{\lstinline!Py\_REFCNT()!} for exact
reference counts. \path{Doc/c-api/refcounting.rst}
warns that immortal objects may report artificial counts and that some
values do not prove uniqueness. In a free-threaded build, a refcount of
one is insufficient evidence that no other thread can access an object.

\section{Ownership ledger}\label{ownership-ledger}

\begin{longtable}[]{@{}
  >{\raggedright\arraybackslash}p{(\columnwidth - 4\tabcolsep) * \real{0.3333}}
  >{\raggedright\arraybackslash}p{(\columnwidth - 4\tabcolsep) * \real{0.3333}}
  >{\raggedright\arraybackslash}p{(\columnwidth - 4\tabcolsep) * \real{0.3333}}@{}}
\toprule\noalign{}
\begin{minipage}[b]{\linewidth}\raggedright
Action
\end{minipage} & \begin{minipage}[b]{\linewidth}\raggedright
Reference effect
\end{minipage} & \begin{minipage}[b]{\linewidth}\raggedright
Code-review question
\end{minipage} \\
\midrule\noalign{}
\endhead
\bottomrule\noalign{}
\endlastfoot
\passthrough{\lstinline!Py\_NewRef(o)!} & Return a new strong reference
& Who decrements it? \\
\passthrough{\lstinline!Py\_INCREF(o)!} & Add a strong reference & Is it
balanced on every exit? \\
\passthrough{\lstinline!Py\_DECREF(o)!} & Release a strong reference &
Could this trigger arbitrary deallocation? \\
\path{Py_XDECREF(o)} & Like DECREF but tolerates
null & Is null an expected state? \\
Borrowed-return API & No new ownership transferred & Can owner disappear
before use? \\
Stealing-reference API & Callee takes ownership on specified conditions
& Does error handling avoid double-free? \\
\end{longtable}

\textbf{ENGINEERING:} Annotate owned/borrowed outputs for every native
helper. When paths branch, draw a ledger for success, failure and
exception exits. Never assume a borrowed pointer survives an operation
that can release the owner's final reference.

\chapter{Cyclic garbage collection}\label{cyclic-garbage-collection}

Reference counting alone cannot collect every unreachable reference
cycle. The cyclic collector detects and disposes of groups of
unreachable objects according to its documented rules.
\passthrough{\lstinline!gc!} exposes control and introspection hooks;
debug results depend on run timing and build configuration. Finalizers
and weak references require care because object resurrection and cleanup
ordering complicate object-graph reasoning.

\begin{lstlisting}[language=Python]
import gc

class Node:
    def __init__(self):
        self.peer = None

left = Node()
right = Node()
left.peer = right
right.peer = left
left = right = None
unreachable = gc.collect()
print("collector-reported unreachable:", unreachable)
\end{lstlisting}

The exact count printed is \emph{not} part of this example's success
criteria. The experiment demonstrates that cycle collection may be
needed; it should not be used to predict object destructor timing in
general.

\section{Weak references and
finalization}\label{weak-references-and-finalization}

A weak reference observes an object without extending its lifetime. Use
\passthrough{\lstinline!weakref!} only for types that support it, and
avoid treating weakref callbacks as an orderly resource shutdown
mechanism. For sockets, files and locks, implement explicit
close/context-management semantics that are robust to delayed GC and
exceptions.

\chapter{Heap allocation, allocator layers and memory
growth}\label{heap-allocation-allocator-layers-and-memory-growth}

Python objects are not identical to operating-system allocations.
CPython may retain arenas, freelists or pools for reuse; freed objects
need not cause immediate decline in process RSS. Analyze memory use
using multiple complementary views: allocation traces
(\passthrough{\lstinline!tracemalloc!}), object graph/GC investigations,
allocator diagnostics and OS process memory. Only compare measurements
under identical interpreter build flags and workload conditions.

\begin{lstlisting}[language=Python]
import tracemalloc

tracemalloc.start()
before = tracemalloc.take_snapshot()
values = [{"key": i} for i in range(10000)]
after = tracemalloc.take_snapshot()
for diff in after.compare_to(before, "lineno")[:5]:
    print(diff)
\end{lstlisting}

This measures tracked Python allocation deltas. It is not a complete
accounting of native allocations inside arbitrary third-party libraries
or allocator-reserved pages.

\chapter{The GIL and concurrency
configurations}\label{the-gil-and-concurrency-configurations}

\section{The GIL is not a data-race correctness
proof}\label{the-gil-is-not-a-data-race-correctness-proof}

Historically, mainstream CPython builds include a global interpreter
lock (GIL) that governs much Python bytecode execution. It does not mean
every compound operation across shared state is atomic, nor does it
replace synchronization around logical invariants. C extensions can
release the GIL and concurrently access external resources. Document
both what a program assumes and what its build actually permits.

\begin{lstlisting}[language=Python]
from threading import Lock

class Counter:
    def __init__(self):
        self._value = 0
        self._lock = Lock()

    def increment(self):
        with self._lock:
            self._value += 1

    def snapshot(self):
        with self._lock:
            return self._value
\end{lstlisting}

\textbf{ENGINEERING:} Evaluate a shared-state algorithm under thread
interleavings independent of whether a particular build happens to
serialize sections of Python execution. For distributed or multiprocess
designs, use explicit IPC and transaction boundaries as well.

\section{Free-threading support}\label{free-threading-support}

\path{Doc/howto/free-threading-python.rst} says
CPython has supported a build with the GIL disabled starting in 3.13.
Free-threaded execution enables multiple cores for parallel threads, but
not every program gets a speedup: lock contention, scheduling, allocator
costs, workload shape and extension behavior remain influential. Some
extension modules may request re-enabling the GIL at import time if they
have not declared support.

The C API provides the \path{Py_GIL_DISABLED}
compile-time macro.
\path{Doc/howto/free-threading-extensions.rst}
describes indicating free-threaded compatibility and adapting C
extensions to thread safety. Declaring support is an engineering
responsibility, not a performance label.

\begin{longtable}[]{@{}
  >{\raggedright\arraybackslash}p{(\columnwidth - 4\tabcolsep) * \real{0.3333}}
  >{\raggedright\arraybackslash}p{(\columnwidth - 4\tabcolsep) * \real{0.3333}}
  >{\raggedright\arraybackslash}p{(\columnwidth - 4\tabcolsep) * \real{0.3333}}@{}}
\toprule\noalign{}
\begin{minipage}[b]{\linewidth}\raggedright
Risk
\end{minipage} & \begin{minipage}[b]{\linewidth}\raggedright
GIL-enabled build
\end{minipage} & \begin{minipage}[b]{\linewidth}\raggedright
Free-threaded build
\end{minipage} \\
\midrule\noalign{}
\endhead
\bottomrule\noalign{}
\endlastfoot
Logical race on shared data & Still possible around multi-step
invariants & Still possible, potentially more concurrency \\
Unsafe native global state & Dangerous & Dangerous and often more
exposed \\
Parallel Python threads & Often limited by GIL & Can execute in parallel
where compatible \\
C extension interaction & GIL policy matters & Explicit compatibility
and synchronization needed \\
Refcount as uniqueness test & Not a reliable general design &
Specifically insufficient for uniqueness \\
\end{longtable}

\chapter{Quiescent-State Reclamation
(QSBR)}\label{quiescent-state-reclamation-qsbr}

\path{InternalDocs/qsbr.md} explains QSBR as a safe
memory-reclamation scheme used by the free-threaded runtime. A logical
removal does not necessarily mean an underlying memory region can be
immediately freed while another thread could still observe it. Threads
periodically report \textbf{quiescent states}, enabling deferred
reclamation once readers are past relevant critical intervals. The
internal document describes applications to dictionary-key arrays, list
backing arrays and allocator pages.

\textbf{Conceptual event timeline:}

\begin{lstlisting}
Thread A                       Thread B
--------                       --------
reads old list array           resizes the list
access may still be live       retires old array
reaches quiescent state        waits for safe reclamation
                               eventually frees old array
\end{lstlisting}

This is a \textbf{teaching abstraction}. It does not specify exact
reclamation ordering or grant extension authors a public QSBR API.
\path{InternalDocs/qsbr.md} is a maintainer
explanation, not a stable concurrency contract for application
developers.

\chapter{Lock ordering and contention
design}\label{lock-ordering-and-contention-design}

\textbf{ENGINEERING controls:} define global lock ordering, keep
blocking I/O out of locks unless deliberate, prefer message passing to
unrestricted shared mutation, and bound work queues to prevent memory
exhaustion. Test for deadlock, liveness and cancellation separately. A
low-contention benchmark can pass even when a rare interleaving corrupts
state.

\textbf{Concurrency invariant template:}

\begin{longtable}[]{@{}
  >{\raggedright\arraybackslash}p{(\columnwidth - 6\tabcolsep) * \real{0.2500}}
  >{\raggedright\arraybackslash}p{(\columnwidth - 6\tabcolsep) * \real{0.2500}}
  >{\raggedright\arraybackslash}p{(\columnwidth - 6\tabcolsep) * \real{0.2500}}
  >{\raggedright\arraybackslash}p{(\columnwidth - 6\tabcolsep) * \real{0.2500}}@{}}
\toprule\noalign{}
\begin{minipage}[b]{\linewidth}\raggedright
Invariant
\end{minipage} & \begin{minipage}[b]{\linewidth}\raggedright
Owner
\end{minipage} & \begin{minipage}[b]{\linewidth}\raggedright
Protected by
\end{minipage} & \begin{minipage}[b]{\linewidth}\raggedright
Failure-injection test
\end{minipage} \\
\midrule\noalign{}
\endhead
\bottomrule\noalign{}
\endlastfoot
Queue length is bounded & Producer/consumer & Semaphore or bounded queue
& Producers saturate the queue \\
Counter increments are not lost & Shared counter & Explicit lock &
High-thread-count increments \\
Cleanup finishes on cancellation & Task worker &
\passthrough{\lstinline!finally!} / context manager & Cancel at each
await point \\
Native module state isolated & Extension & Per-module/interpreter design
& Multiple subinterpreters \\
\end{longtable}

\chapter{Native concurrency review}\label{native-concurrency-review}

\path{Doc/c-api/threads.rst},
\path{Doc/c-api/synchronization.rst},
\path{Doc/c-api/subinterpreters.rst},
\path{Doc/howto/free-threading-extensions.rst}, and
\path{Include/critical_section.h} cover relevant
facets, at varying stability levels. The presence of functions or
internal headers is not itself a promise that an extension should access
them directly. Choose public APIs according to the declared API
stability tier.

\chapter{Laboratories and evidence}\label{laboratories-and-evidence}

\begin{enumerate}
\def\labelenumi{\arabic{enumi}.}
\tightlist
\item
  \textbf{Cycle and allocation experiment:} build a cyclic object graph,
  collect it, measure allocations using
  \passthrough{\lstinline!tracemalloc!}, and explain why RSS can remain
  high.
\item
  \textbf{Thread-safety experiment:} run a shared-counter workload with
  and without deliberate synchronization; do not rely on reproducing a
  race to prove absence of a race.
\item
  \textbf{Lifetime audit:} review a C helper that returns borrowed and
  owned references. Mark the ownership at every return/exception branch.
\item
  \textbf{Free-threading compatibility plan:} inventory external C
  extension modules, their declared GIL support and uses of
  global/static mutable state; create a matrix of blockers.
\item
  \textbf{QSBR design explanation:} draw a two-thread sequence diagram
  for resizing a shared container and delayed memory reclamation; mark
  it explicitly as an internal implementation model.
\end{enumerate}

\textbf{Completion rubric:} source citation accuracy (30\%), correct
lifetime analysis (30\%), reproducible tests (25\%), and precise
separation of public guarantees from internal mechanics (15\%).

\chapter{Source concordance}\label{source-concordance}

\begin{itemize}
\tightlist
\item
  \textbf{SOURCE:}
  \path{InternalDocs/garbage_collector.md},
  \passthrough{\lstinline!qsbr.md!},
  \passthrough{\lstinline!stackrefs.md!},
  \passthrough{\lstinline!frames.md!},
  \path{interpreter.md}.
\item
  \textbf{SOURCE:} \path{Doc/c-api/refcounting.rst},
  \passthrough{\lstinline!memory.rst!},
  \path{gcsupport.rst},
  \passthrough{\lstinline!threads.rst!},
  \path{synchronization.rst},
  \path{subinterpreters.rst}.
\item
  \textbf{SOURCE:}
  \path{Doc/howto/free-threading-python.rst},
  \path{free-threading-extensions.rst},
  \path{Doc/library/gc.rst},
  \path{tracemalloc.rst}.
\item
  \textbf{IMPLEMENTATION:} \path{Python/qsbr.c},
  \path{Python/gc_gil.c},
  \passthrough{\lstinline!Python/brc.c!},
  \path{Include/internal/pycore_qsbr.h} as present
  in the snapshot body inventory.
\item
  \textbf{ENGINEERING:} all design invariants, synchronization examples,
  illustrative timelines and labs.
\end{itemize}

\end{document}
